\noindent\textbf{Elección del procedimiento.} La exigencia es $a_c=ng$, no una fuerza $n$ veces la masa. En giro uniforme $a_c=\Omega^2R$; se despeja la raíz positiva de $\Omega$. El efecto sobre la cabeza es un segundo cálculo, con su velocidad relativa radial.

\begin{enumerate}

\item De $\Omega^2R=ng$, dividir por $R$ y tomar la raíz.

\[\fitmath{\Omega=\sqrt{\frac{n\cdot g}{R}}}\]

\[\fitmath{\Omega=\sqrt{\frac{(5)\cdot (9.81)}{(8)}}\approx 2.476136\,\mathrm{1/s}}\]

\item Una vuelta son $2\pi$ radianes y un minuto son 60 s: convertir rad/s a vueltas/min.

\[\fitmath{N=\frac{\Omega\cdot 60}{2\cdot \pi}}\]

\[\fitmath{N=\frac{(2.476136103)\cdot 60}{2\cdot \pi}\approx 23.64536\,\mathrm{min^{-1}}}\]

\item La velocidad radial es perpendicular a la rotación, de modo que el módulo de $2\boldsymbol\Omega\times\boldsymbol v$ es $2\Omega v$.

\[\fitmath{\mathrm{aC}=2\cdot \Omega\cdot v}\]

\[\fitmath{\mathrm{aC}=2\cdot (2.476136103)\cdot (0.5)\approx 2.476136\,\mathrm{m/s^{2}}}\]

\end{enumerate}

\noindent\textbf{Contraste y comparación.} El vector de Coriolis es transversal al movimiento radial; su sentido depende de si la cabeza se mueve hacia dentro o hacia fuera. Verifica $\Omega^2R/g=5$.
